\section{Conclusion et Perspectives}

\noindent Cet article a présenté la conception, l'implémentation et l'évaluation
rigoureuse d'un système de Question-Réponse juridique marocain basé sur
l'architecture Retrieval-Augmented Generation (RAG). Le système intègre une
chaîne complète de traitement du langage naturel arabe — extraction PDF,
nettoyage, normalisation arabe, analyse structurelle, indexation vectorielle,
recherche hybride, et génération de réponses par LLM — dans une plateforme
open-source modulaire et reproductible.

\vspace{0.2cm}

\noindent Les contributions principales de ce travail sont :

\begin{enumerate}
    \item \textbf{Une chaîne de NLP arabe complète et robuste} pour le
    domaine juridique marocain, incluant un extracteur PDF multi-méthode
    (pdfplumber, PyMuPDF, EasyOCR), un nettoyeur de bruit spécifique aux
    textes juridiques, un normaliseur arabe standardisé, et un analyseur
    structurel de codes de loi.
    \item \textbf{Une architecture de recherche hybride dense + BM25 + RRF}
    indexée dans PostgreSQL + pgvector avec index HNSW, démontrant que la
    fusion par rangs améliore le Hit@5 à 68,75\% et le Recall@5 à 59,38\%.
    \item \textbf{Le premier benchmark comparatif de cinq modèles d'embedding
    multilingues} sur un corpus juridique arabe avec 65 variantes de requêtes,
    établissant la supériorité des modèles 1024 dimensions (E5-large : MRR
    0,81, précision de domaine 89\%).
    \item \textbf{Une évaluation complète de la qualité des réponses générées}
    par Gemma 4 2B, révélant une fidélité de 22\% comme axe d'amélioration
    prioritaire.
\end{enumerate}

\vspace{0.2cm}

\noindent Le tableau~\ref{tab:comparaison} positionne notre système par rapport
aux alternatives existantes pour le droit marocain.

\begin{table}[H]
\centering
\caption{Comparaison avec les solutions existantes}
\label{tab:comparaison}
\begin{tabular}{lccccc}
\toprule
\textbf{Capacité} & \textbf{Google} & \textbf{legislation.ma} & \textbf{ChatGPT} & \textbf{LLM seul} & \textbf{Notre système} \\
\midrule
Recherche sémantique arabe & Non & Non & Non & Non & \textbf{Oui} \\
Réponses générées & Non & Non & Oui & Oui & \textbf{Oui} \\
Citations d'articles & -- & -- & Parfois & Parfois & \textbf{Oui} \\
Ancrage dans les sources & -- & -- & Non & Non & \textbf{Oui} \\
Open-source & -- & -- & Non & Variable & \textbf{Oui} \\
Déploiement local (Docker) & -- & -- & Non & Variable & \textbf{Oui} \\
Corpus juridique marocain & Partiel & Oui & Non & Non & \textbf{Oui} (7 domaines) \\
\bottomrule
\end{tabular}
\end{table}

\noindent Les résultats confirment la viabilité de l'approche RAG pour le
question-réponse juridique arabe. Les trois campagnes de benchmark ont
démontré : (1) la supériorité des modèles d'embedding 1024 dimensions
(E5-large : MRR 0,81, précision de domaine 89\%) ; (2) l'efficacité de la
recherche hybride Dense + RRF (Hit@5 = 68,8\%, Recall@5 = 59,4\%) ; (3) la
faisabilité d'un déploiement léger via un LLM quantifié de $\sim$1,1 Go.

\vspace{0.2cm}

\noindent Les limites identifiées — gap dialectal Darija--MSA (29\%), latence
du pipeline (13--14~s), couverture partielle du corpus, domaine fiscal
inaccessible — définissent une feuille de route claire : fine-tuning
embedding sur corpus juridique, déploiement local de l'embedding, et
évaluation par LLM-as-judge pour une mesure plus juste de la fidélité.

\vspace{0.2cm}

\noindent Au-delà de ses résultats techniques, ce projet démontre le potentiel
de l'architecture RAG pour démocratiser l'accès au droit dans les pays où la
barrière linguistique et la dispersion des sources constituent un obstacle
majeur. La méthodologie développée — extraction OCR+VLM, normalisation arabe,
embedding multilingue, recherche hybride, génération ancrée — est transposable
à d'autres corpus juridiques arabophones (Tunisie, Algérie, Égypte) avec un
effort de localisation minimal.